\documentclass[a4paper]{article}
\usepackage[utf8]{inputenc}
\usepackage[T2A]{fontenc}
\usepackage[english, russian]{babel}
\usepackage[
    left = 25mm,
    top = 20mm,
    right = 25mm,
    bottom = 30mm,
    headheight = 14pt,
    footskip = 15pt
]{geometry}

\usepackage{amsmath, amsfonts, amssymb, amsthm} 
\usepackage{fancyhdr, fancybox}
\usepackage{graphicx}
\usepackage{xcolor}
\usepackage{hyperref}
\usepackage{tikz}
\usepackage{float}

\usetikzlibrary{arrows.meta}
\hypersetup{colorlinks = true, allcolors = [RGB]{010 090 200}}
\renewcommand{\proofname}{\textbf{Решение}}
\newcommand{\BigTau}{\scalebox{1.4}{$\tau$}}
\pagestyle{fancy}
\fancyhf{}
\fancyfoot[R]{\thepage}
\fancyhead[L]{Наглядная топология и геометрия}
\fancyhead[R]{\href{https://github.com}{Барменков Федор}}
\headsep = 10 mm

\begin{document}

\section*{Домашнее задание для допуска на экзамен}

\subsection*{Упражнение 1}
\begin{proof}
Предположим, у всех людей в компании различное количество знакомых. Тогда количество знакомых может быть от \ $0$ \  до \ $n - 1$. В таком случае в компании одновременно существует человек, знакомый со всеми и человек без знакомых. Такого быть не может - противоречие.\end{proof}

\subsection*{Упражнение 2.1}
\begin{proof}
Топология Зарисского на множестве $X$ задается семейством открытых подмножеств:
\[\BigTau_{\text{З}} = \left\{ U \subset X \;\middle|\; \lvert X \setminus U \rvert < \infty \right\} \cup \{\varnothing\}\]

Проверим, что $\BigTau_{\text{З}}$ действительно удовлетворяет аксиомам топологической структуры:

\begin{enumerate}
    \item Аксиома пустого множества и всего пространства:
    
    По определению, $\varnothing \in \BigTau_{\text{З}}$. 
    Так же $X \in \BigTau_{\text{З}}$, поскольку дополнение всего пространства до самого себя пусто:
    \[\lvert X \setminus X \rvert = \lvert \varnothing \rvert = 0 < \infty\]
    
    \item Аксиома пересечения:
    
    Пусть $U, V \in \BigTau_{\text{З}}$. Докажем, что их пересечение $U \cap V \in \BigTau_{\text{З}}$.
    Если $U \cap V = \varnothing$, то очевидно.
    Предположим, что $U \cap V \neq \varnothing$, тогда: \[X \setminus (U \cap V) = (X \setminus U) \cup (X \setminus V)\]
    Так как $\lvert X \setminus U \rvert < \infty$ и $\lvert X \setminus V \rvert < \infty$, то и $\lvert X \setminus (U \cap V) \rvert < \infty$, а значит $U \cap V \in \BigTau_{\text{З}}$.
    Продолжив по индукции, выстроим конечный набор подмножеств $\{U_\alpha\}_{\alpha \in I}$ имеем: \[
    X \setminus \left( \bigcap_{\alpha \in I} U_{\alpha} \right) = \bigcup_{\alpha \in I} (X \setminus U_{\alpha}). \ \text{Так как} \ \lvert X \setminus U_\alpha \rvert < \infty \implies\lvert \bigcup_{\alpha \in I} (X \setminus U_{\alpha}) \rvert < \infty \implies \bigcap_{\alpha \in I} U_{\alpha} \in \BigTau_{\text{З}}\]

    \item 

    Аксиома объединения:
    
    Пусть $\{U_\alpha\}_{\alpha \in I} \subset \BigTau_{\text{З}}$. Причем $\lvert I \rvert \ge \aleph_0$. Докажем, что их 
    $\bigcup_{\alpha \in I} U_\alpha \in \BigTau_{\text{З}}. \ \text{Если все} \ U_\alpha = \varnothing,\ \text{то очевидно}.$
    Если хотя бы один $U_{\alpha_0}$ не пустой, то дополнение к нему  $\lvert X \setminus U_{\alpha_0} \rvert$ конечно. По закону де Моргана имеем: \[X \setminus \left( \bigcup_{\alpha \in I} U_\alpha \right) = \bigcap_{\alpha \in I} X \setminus U_\alpha \ \text{и} \ \bigcap_{\alpha \in I} X \setminus U_\alpha \subset X \setminus U_{\alpha_0} \implies \lvert \bigcap_{\alpha \in I} \left(X \setminus U_\alpha \right) \rvert < \infty \implies \bigcup_{\alpha \in I} U_\alpha \in \BigTau_{\text{З}}\]
\end{enumerate}
\end{proof}
\subsection*{Упражнение 2.2}
\begin{proof}
Пусть $(X, \BigTau_X)$ - топологическое пространство и $Y \subset X$, тогда индуцированная топология задается пересечением подмножеств $X$ и самого множества $Y$:
\[\BigTau_Y = \left\{U \cap Y  \ | \ U \in \BigTau_X\right\}\]
\begin{enumerate}
    \item Аксиома пустого множества и всего пространства:

    Так как $\varnothing \in \BigTau_X \implies \varnothing \cap Y = \varnothing \in \BigTau_Y$ и $Y \subset X \implies Y \cap Y = Y \in \BigTau_Y$ 

    \item Аксиома пересечения:

    Пусть $U, V \in \BigTau_Y$. Докажем, что их пересечение $U \cap V \in \BigTau_Y$. Если $U \cap V = \varnothing$, то очевидно. Иначе, по закону де Моргана:
    \[(U \cap Y) \cap (V \cap Y) = (U \cap V) \cap Y. \ \text{Так как} \ \BigTau_X \ \text{- топология, то} \ U \cap V \in \BigTau_X \implies U \cap V \in \BigTau_Y\]

    \item Аксиома объединения: 

    Пусть $\{U_\alpha\}_{\alpha \in I} \subset \BigTau_Y$ - набор открытых множеств. Докажем, что: \[\bigcup_{\alpha \in I} U_\alpha \in \BigTau_Y\]
    По аналогии с 2 пунктом упражнения имеем: \[\bigcup_{\alpha \in I}(U_\alpha \cap Y) = Y \cap \bigcup_{\alpha \in I} U_\alpha \implies \bigcup_{\alpha \in I} U_\alpha \in \BigTau_Y\]
\end{enumerate}
\end{proof}
\subsection*{Упражнение 3.1}
\begin{proof}
Рассмотрим топологическое пространство $(X, \BigTau_X)$ с антидискретной топологией, где:
\[
\BigTau_X = \{ \varnothing, X \}
\]
Рассмотрим произвольную последовательность $\left\{ x_n \right\}_{n=1}^{\infty}$ и предположим, что она имеет предел $A \in X$. По определению предела:
\[\forall \, U_A \in \BigTau_X \quad \exists \, N \in \mathbb{N} : \quad \forall \, n \ge N \quad x_n \in U_A\]
Поскольку из окрестностей точки $A$ у нас выбор довольно скромный - само множество $X$, то для любой последовательности нам подойдет любое число из исходного множества, а значит наша последовательность сходится ко всем точкам множества $X$.
\end{proof}

\subsection*{Упражнение 3.3}
\begin{proof}
Напомним, что топологическое пространство называется хаусдорфовым, если для любых двух различных точек $x, y$ существуют такие их окрестности $U_x$ и $V_y$, что:
\[U_x \cap V_y = \varnothing\]
Рассмотрим последовательность $\left\{ x_n \right\}_{n=1}^{\infty}$. Предположим от противного, что она имеет два различных предела $A$ и $B$ ($A \ne B$). Тогда по определению предела для любых их окрестностей $U_A$ и $V_B$ должны существовать номера $N_1, N_2 \in \mathbb{N}$, такие что:\[
\forall \, n \ge N_1 \implies x_n \in U_A \quad \text{и} \quad \forall \, n \ge N_2 \implies x_n \in V_B\]
Положив $N = \max\left\{N_1, N_2\right\}$, для всех $n \ge N$ мы получаем:
\[x_n \in U_A \cap V_B \implies U_A \cap V_V \ne \varnothing\]
А поскольку окрестности выбирались произвольно - это противоречит определению хаусдорфовости. Значит, предел хаусдорфова пространства единственен.
\end{proof}
\subsection*{Упражнение 4.2}
\begin{proof}
Рассмотрим две последовательности на числовой прямой:
\[x_n = n \quad \text{и} \quad y_n = n + \frac{1}{n} \quad \forall n \in \mathbb{N}\]
\begin{figure}[H]
    \centering
    \begin{tikzpicture}[
        scale=2.3,
        every node/.style={font=\small},
        axis/.style={thick, ->, >=stealth},
        seq1/.style={circle, fill=blue, inner sep=2pt},
        seq2/.style={circle, fill=red, inner sep=2pt}
    ]
        \draw[axis] (-0.05,0) -- (4.8,0) node[right] {$\mathbb{R}$};
        
        \draw[thick] (0, 0.05) -- (0, -0.05) node[below=2pt] {0};
        \draw[thick] (1, 0.05) -- (1, -0.05) node[below=2pt] {1};
        \draw[thick] (2, 0.05) -- (2, -0.05) node[below=8pt] {2};
        \draw[thick] (3, 0.05) -- (3, -0.05) node[below=2pt] {3};
        \draw[thick] (4, 0.05) -- (4, -0.05) node[below=2pt] {4};

        \node[seq1, label=above:{\color{blue}$x_1$}] at (1, 0) {};
        \node[seq1, label=above:{\color{blue}$x_2$}] at (2, 0) {};
        \node[seq1, label=above:{\color{blue}$x_3$}] at (3, 0) {};
        \node[seq1, label=above:{\color{blue}$x_4$}] at (4, 0) {};
        \node[blue] at (4.5, 0.15) {$\dots$};

        \node[seq2, label=below:{\color{red}$y_1$}] at (2, 0) {};
        \node[seq2, label=below:{\color{red}$y_2$}] at (2.5, 0) {};
        \node[seq2, label=below:{\color{red}$y_3$}] at (3.333, 0) {};
        \node[seq2, label=below:{\color{red}$y_4$}] at (4.25, 0) {};
        \node[red] at (4.6, -0.2) {$\dots$};

        \node[anchor=west, blue] at (0, 0.5)  {$\bullet \ x_n = n$};
        \node[anchor=west, red]  at (0, 0.3) {$\bullet \ y_n = n + \frac{1}{n}$};
    \end{tikzpicture}
    \caption{Иллюстрация последовательностей $x_n$ и $y_n$ на числовой оси}
\end{figure}
Понятно, что наши числовые последовательности:
\[X = \left\{x_n\right\}_{n=1}^{\infty} \quad \text{и} \quad Y = \left\{y_n\right\}_{y=1}^{\infty}\]
Образуют два непересекающихся и замкнутых подмножества на числовой прямой. Вспомним что расстояние между множествами определяется как:
\[\rho(X, Y) = \inf \left\{\lvert x - y \rvert \;\middle|\; x \in X, y \in Y\right\} \implies \rho(X, Y) = \inf \left\{ \frac{1}{n} \;\middle|\; n \in \mathbb{N} \right\} = \lim_{n \rightarrow \infty} \frac{1}{n} = 0\] 
Значит это два непустых непересекающихся подмножества в $\mathbb{R}$, расстояние между которыми ноль.
\end{proof}

\end{document}
